\chapter{Linear Dimension Reduction via Random Projections}
\label{c:johnson}



In Chapters~\ref{c:svd} and~\ref{c:diffusion} we saw both linear and non-linear methods for dimension reduction.
In this chapter we will see one of the most fascinating consequences of the phenomenon of concentration of measure in high dimensions, one of the \emph{blessings of high dimensions} described in Chapter~\ref{c:surprises}. When given a data set in high dimensions, we will see  that the projection to a lower dimensional space, taken at random, can preserve (under conditions made precise in the next section) certain geometric features of the data set. The remarkable aspect here is that this ``lower'' dimension can be strikingly lower. This gives rise to significant computational savings in many data processing tasks by including a random projection as a pre-processing step. Similar ideas are also the driving force behind many algorithms in {\em randomized linear algebra}, as we will demonstrate in this chapter via the {\em randomized SVD}.
There is however another less obvious implication of this phenomenon with important practical implications: since the projection is agnostic of the data, it can be leveraged even when the data set is not explicit, such as the set of all \emph{natural images} or the set of all ``possible'' \emph{brain scans}; this is at the heart of \emph{Compressive Sensing}, which we will explore in Chapter~\ref{c:cs}



\section{The Johnson-Lindenstrauss Lemma}\label{s:jl}

Suppose one has $n$ points, $X = \{x_1,\dots,x_n\}$, in $\RR^p$ (with $p$ large). If $p>n$, the points actually lie in a subspace of dimension $n$, so the projection $f:\RR^p\to \RR^n$ of the points to that subspace acts without distorting the geometry of $X$. In particular, for every $x_i$ and $x_j$, $\|f(x_i)-f(x_j)\|^2 = \|x_i-x_j\|^2$, meaning that $f$ is an isometry in $X$.
Suppose instead we allow a bit of distortion, and look for a map $f:\RR^p\to \RR^d$ that is an $\epsilon-$isometry\footnote{Another possible convention is to ask for $(1-\epsilon) \|x_i-x_j\| \leq \|f(x_i)-f(x_j)\| \leq (1+\epsilon) \|x_i-x_j\|$; this is equivalent up to a multiplicative constant on $\epsilon$ since for $0\leq\epsilon\leq 1$ we have $1+\epsilon\leq (1+\epsilon)^2\leq 1+3\epsilon$ and $1-2\epsilon\leq (1-\epsilon)^2\leq 1-\epsilon$.}, meaning that
\begin{equation}\label{def_epsilonisometry}
(1-\epsilon) \|x_i-x_j\|^2 \leq \|f(x_i)-f(x_j)\|^2 \leq (1+\epsilon) \|x_i-x_j\|^2.
\end{equation}
Can we do better than $d=n$?

In 1984, Johnson and Lindenstrauss~\cite{Johnson_Lindenstrauss_84} showed a remarkable lemma that answers this question affirmatively.

\begin{theorem}[Johnson-Lindenstrauss Lemma~\cite{Johnson_Lindenstrauss_84}]\label{JL_random_0}
 For any $0<\epsilon<1$ and for any integer $n$, let $d$ be such that
 \begin{equation}\label{kdim}
 d \geq 4\frac1{\epsilon^2/2 - \epsilon^3/3}\log n.
\end{equation}
Then, for any set $X$ of $n$ points in $\RR^p$, there is a linear map $f:\RR^p\to\RR^d$ that is an $\epsilon-$isometry for $X$ (see (\ref{def_epsilonisometry})).
This map can be found in randomized polynomial time\footnote{Randomized polynomial time refers to the complexity class of decision problems for which a randomized algorithm can solve the problem in polynomial time with a high probability of correctness.}.
\end{theorem}



\begin{proof}
We follow~\cite{Dasgupta_Gupta_JLsimpleproof} for an elementary proof for Theorem~\ref{JL_random_0}. 
We will start by showing that, given a pair $x_i,x_j$ a projection onto a random subspace of dimension $d$ will satisfy (after appropriate scaling) property (\ref{def_epsilonisometry}) with high probability. Without loss of generality we can assume that $u = x_i - x_j$ has unit norm. Understanding what is the norm of the projection of $u$ on a random subspace of dimension $d$ is the same as understanding the norm of the projection of a (uniformly) random point on $S^{p-1}$ the unit sphere in $\RR^p$ on a specific $d$-dimensional subspace---let us say the one generated by the first $d$ canonical basis vectors.

This means that we are interested in the distribution of the norm of the first $d$ entries of a random vector drawn from the uniform distribution over $S^{p-1}$ -- this distribution is the same as taking a standard Gaussian vector in $\RR^p$ and normalizing it to the unit sphere.

Let $g:\RR^p\to\RR^d$ be the projection on a random $d-$dimensional subspace and let $f:\RR^p\to\RR^d$ defined as $f = \sqrt{\frac{p}d}g$. Then (by the above discussion), given a pair of distinct points $x_i$ and $x_j$, $\frac{\|f(x_i)-f(x_j)\|^2}{\|x_i-x_j\|^2}$ has the same distribution as $\frac{p}dL$, as defined in Lemma~\ref{lemma_COM}. Using Lemma~\ref{lemma_COM}, we have, given a pair $x_i,x_j$,

\[
 \Pr\left[ \frac{\|f(x_i)-f(x_j)\|^2}{\|x_i-x_j\|^2} \leq (1-\epsilon) \right] \leq \exp\left( \frac{d}2(1-(1-\epsilon) +\log(1-\epsilon)) \right),
\]
since for $\epsilon\geq 0$, $\log(1-\epsilon) \leq -\epsilon -\epsilon^2/2$, using~\eqref{kdim} we have
\begin{eqnarray*}
 \Pr\left[ \frac{\|f(x_i)-f(x_j)\|^2}{\|x_i-x_j\|^2} \leq (1-\epsilon) \right] & \leq & \exp\left( -\frac{d\epsilon^2}4 \right) \\
      & \leq & \exp\left( -2\log n \right) = \frac1{n^2}.
\end{eqnarray*}
On the other hand,
\[
 \Pr\left[ \frac{\|f(x_i)-f(x_j)\|^2}{\|x_i-x_j\|^2} \geq (1+\epsilon) \right] \leq \exp\left( \frac{d}2(1-(1+\epsilon) +\log(1+\epsilon)) \right).
\]
since for $\epsilon\geq 0$, $\log(1+\epsilon) \leq \epsilon -\epsilon^2/2 + \epsilon^3/3$, using~\eqref{kdim} we have
\begin{eqnarray*}
 \Prob\left[ \frac{\|f(x_i)-f(x_j)\|^2}{\|x_i-x_j\|^2} \geq (1+\epsilon) \right] & \leq & \exp\left( -\frac{d\left(\epsilon^2 - 2\epsilon^3/3\right)}4 \right) \\
      & \leq & \exp\left( -2\log n \right) = \frac1{n^2}.
\end{eqnarray*}
By the union bound it follows that 
\[
 \Pr\left[ \frac{\|f(x_i)-f(x_j)\|^2}{\|x_i-x_j\|^2} \notin [1-\epsilon,1+\epsilon] \right] \leq \frac2{n^2}.
\]
Since there exist ${n \choose 2}$ such pairs, again, a simple union bound gives
\[
 \Pr\left[ \exists_{i,j}:\frac{\|f(x_i)-f(x_j)\|^2}{\|x_i-x_j\|^2} \notin [1-\epsilon,1+\epsilon] \right] \leq \frac2{n^2} \frac{n(n-1)}{2} = 1-\frac1n.
\]
Therefore, choosing $f$ as a properly scaled projection onto a random $d$-dimensional subspace gives an $\epsilon$-isometry on $X$ (see (\ref{def_epsilonisometry})) with probability at least $\frac1n$. We can achieve any desirable constant probability of success by trying $\OOO(n)$ such random projections, meaning we can find an $\epsilon-$isometry in randomized polynomial time.

\end{proof}

Note that by considering $d$ slightly larger one can get a good projection on the first random attempt with  high confidence. In fact, it is trivial to adapt the proof above to obtain the following lemma:

\begin{proposition}\label{JL_random}
 For any $0<\epsilon<1$, $\tau>0$, and for any integer $n$, let $d$ be such that
\[
 d \geq \frac{2(2+\tau)}{\epsilon^2/2 - \epsilon^3/3}\log n.
\]
Then, for any set $X$ of $n$ points in $\RR^p$, take $f:\RR^p\to\RR^d$ to be a suitably scaled projection on a random subspace of dimension $d$, then $f$ is an $\epsilon-$isometry for $X$ (see (\ref{def_epsilonisometry})) with probability at least $1-\frac1{n^{\tau}}$.
\end{proposition}

Proposition \ref{JL_random} is quite remarkable. Consider the situation where we are given a high-dimensional data set in a streaming fashion -- meaning that we get each data point at a time, consecutively. To run a dimension-reduction technique like PCA or Diffusion maps we would need to wait until we received the last data point and then compute the dimension reduction map (both PCA and Diffusion Maps are, in some sense, data adaptive). Using Lemma \ref{JL_random} one can just choose a projection at random in the beginning of the process (all one needs to know is an estimate of the logarithm of the size of the data set) and just map each point using this projection matrix which  can be done online -- we do not need to see the next point to compute the projection of the current data point. Proposition~\ref{JL_random} ensures that this (seemingly na\"ive) procedure will, with high probably, not distort the data by more than $\epsilon$.




One might wonder if such low-dimensional embeddings as provided by the Johson-Lindenstrauss Lemma also extend to other norms besides the Euclidean norm.  For the $\ell_1$-norm there exist lower bounds which prevent such a dramatic dimension reduction (see~\cite{lee2004embedding}), and for the $\ell_{\infty}$-norm one can easily construct examples that demonstrate the impossibility of dimension reduction. Hence, the Johnson-Lindenstrauss Lemma seems to be a subtle result about the Euclidean norm.




\subsection{The Fast Johnson-Lindenstrauss transform}

Let us continue thinking about our example of high-dimensional streaming data. After we draw the random projection matrix\footnote{An orthogonal projection $P$ must satisfy $P=P^{\ast}$ and $P^2=P$. Here, it is not $M$ that represents a projection, but $M^{\ast} M$, yet for our purposes of approximate norm-preserving dimension reduction it suffices to apply $M$ instead of $M^{\ast} M$.  However, with a slight abuse of terminology, we still refer to $M$ as projection.}, say $M$, for each data point $x$, we still have to compute $Mx$ which has a computational cost of $\OOO(\epsilon^{-2}\log(n)p)$  since $M$ has $\OOO(\epsilon^{-2}\log(n)p)$ entries (since $M$ is a random matrix, generically it will be a dense matrix).  In some applications this might be too expensive, raising the natural question of whether one can do better. 
Moreover, storing a large-scale dense matrix $M$  is not very desirable either.
There is no hope of significantly reducing the number of rows in general, as it is known that  the Johnson-Lindenstrauss Lemma is orderwise optimal~\cite{Alon_JL_lowerbound_2003,Larsen_Nelson_JL}. 


We might hope to replace the dense random matrix $M$ by a sparse matrix $S$ to speed up the matrix-vector multiplication and to  reduce the storage requirements.  This method was proposed and analyzed in~\cite{Ailon_Chazelle_fast_JL}. Here we discuss a slightly simplified version,  see also~\cite{cohen2015optimal}.

We let $S$ be a very sparse $d \times p$ matrix, where each row of $S$ has just one single non-zero entry of value
$\sqrt{p/d}$ at a location drawn uniformly at random. Then, for any vector $x \in \R^p$ 
\begin{equation*}
  \underset{i}{\E}
   [(Sx)_i^2]=\sum\limits_{j=1}^p \P(S_{ij}\neq 0) \cdot \frac{p}{d} \cdot x^2_j = \frac{1}{d} \|x\|^2,
\end{equation*}
hence $\E [\|Sx\|^2] = \E [\sum\limits^d_{i = 1} (Sx_i)^2]=\|x\|^2$.
This result is satisfactory with respect to expectation (even for $d=1$), but not with respect to the variance of $\|Sx\|^2$. For instance,
if $x$ has only one non-zero entry  we need $d \sim \OOO(p)$ to ensure that $\|Sx\|^2 \neq 0$ with non-negligible probability. More generally, if one coordinate of $x$ is much larger (in absolute value) than all its other coordinates, then we will need a rather large value for $d$ to guarantee that $\|Sx\| \approx \|x\|$.

A natural way to quantify the ``peakiness'' of a vector $x$ is via the {\em peak-to-average ratio}\footnote{This quantity also plays an important role in wireless communications. There, one tries to avoid transmitting signals with a large peak-to-average ratio, since such signals would suffer from nonlinear distortions when they are passing through  the power amplifiers that are usually installed in cell phones. The potentially large 
peak-to-average ratio of OFDM signals is one of the alleged reasons why CDMA was dominant over OFDM for such a long time.}
measured by the quantity $\|x\|_\infty/\|x\|$.  It is easy to see that we have (assuming $x$ is not the zero-vector)
$$ \frac{1}{\sqrt{p}} \le \frac{\|x\|_\infty}{\|x\|} \le 1.$$
The upper bounds is achieved by vectors with only one non-zero entry, while the lower bound is met by constant-modulus vectors.
Thus, if  
\begin{equation}\label{approxpar}
\frac{\|x\|_\infty}{\|x\|} \approx \frac{1}{\sqrt{p}},
\end{equation}
we can hope that sparse subsampling of $x$ will still preserve its Euclidean norm.

Thus, this suggests to include a preprocessing step by applying a rotation so that sparse vectors become non-sparse in the new basis, thereby reducing their $\infty$-norm (while their 2-norm remains invariant under rotation). Two natural choices for such a rotation are  the Discrete Fourier transform (which maps unit-vectors into constant modulus vectors) and its $\Z_2$-cousin, the Walsh-Hadamard matrix\footnote{Hadamard matrices do not exist for all dimension $d$. But we can always pad $x$ with zeroes to achieve the desired length.}.
But since the Fast Johnson-Lindenstrauss Transform (FJLT) has to work for all vectors, we  need to avoid that this rotation maps dense vectors into sparse vectors. We can address this issue by ``randomizing'' the rotation, thereby ensuring with overwhelming probability that dense vectors are not mapped into sparse vectors. This can be accomplished in a numerically efficient manner (thus maintaining our overall goal of numerical efficiency) by first randomizing the signs of $x$ before applying the rotation. Putting these steps together we arrive at the following map.
\begin{definition}\label{FastJL}
The {\em Fast Johnson-Lindenstrauss Transform} is  the map $\Psi: \C^p \to \C^d$, defined by $\Psi:= S F D$, where 
$S$ and $D$ are random matrices and $F$ is a deterministic matrix. In particular,
\begin{itemize}
\setlength{\itemsep}{-0.5ex}
\setlength{\parsep}{-0.5ex}
\item $S$ is a $d \times p$ matrix, where each row of $S$ has just one single non-zero entry of value $\sqrt{p/d}$ at a location drawn uniformly at random.
\item $F$ is either the $p \times p$ DFT matrix or the $p \times p$ Hadamard matrix (if it exists), in each case normalized by $1/\sqrt{p}$ to obtain a unitary matrix.
\item $D$ is a $p\times p$ diagonal matrix whose entries are drawn independently  from $\{-1,+1\}$ with probability $1/2$.
\end{itemize}
\end{definition}

We can carry out the matrix-vector multiplication by the DFT matrix via the Fast Fourier Transform (FFT) in $\OOO(p \log p)$ operations; a similar algorithm exists for the Walsh-Hadamard matrix.
The FJLT allows for a dimension reduction that is competitive with the Johnson-Lindenstrauss Lemma as manifested by the following theorem.

\begin{theorem}[Fast Johnson-Lindenstrauss Transform]\label{th:fjl}
For $0<\eps<1$ and $0<\delta<1$, there is a random matrix $\Psi$ of size $d \times p$ with $d = \OOO \big( \log(p/\delta) \log(1/\delta) /\eps^2\big)$ such that, for
each $x \in \C^p$,
$$\| \Psi x \| \in [1-\eps,1+\eps] \cdot \|x\|$$
holds with probability at least $1-\delta$. Matrix-vector multiplication with $\Psi$ takes $\OOO(p\log p +d)$ operations.
\end{theorem}

For convenience we will prove Theorem~\ref{th:fjl} for the case when $F$ is the Walsh-Hadamard matrix so that the random variables are real-valued and the exposition is lighter. But it is straighforward to adapt the argument to the case when $F$ is the DFT matrix. Keeping this minor simplification in mind, the proof of Theorem~\ref{th:fjl} follows from the two lemmas below.
We first show that with high probability the random rotation $FD$ produces vectors with a sufficiently low peak-to-average ratio.

\begin{lemma}\label{le:par2}
Let $y = FDx$, where $F$ and $D$ are as in Definition~\ref{FastJL}. Then
\begin{equation}\label{jlpar}
\underset{D}{\P} \left( \frac{\| y \|_\infty}{\|y\|} \ge  \sqrt{\frac{2 \log (4p/\delta)}{p}}\right) \le \frac{\delta}{2}.
\end{equation}
\end{lemma}

\begin{proof}
Since $FD$ is unitary, the quantity $\|FDx\|_\infty/\|FDx\|$ is invariant under rescaling of $x$ and therefore we can assume $\|x\|=1$.

Let $\xi_i = \pm 1$ be the $i$-th diagonal entry of $D$. We have $y_i  =  \sum_{j=1}^p \xi_j F_{ij} x_j$ and note that the terms of this sum are i.i.d.\ bounded random variables. We thus can apply Hoeffding's inequality. In the notation of Theorem~\ref{thm:hoeffding1},
let $X_j =  \xi_j F_{ij} x_j$. We note that $X_j = \pm  F_{ij} x_j$,  hence $\E[X_j] =0$ and 
$|X_j  | \le a_j$, where $a_j = \left|F_{ij} x_j\right|$. It holds that
$$\sum_{j=1}^p a_j^2 = \sum_{j=1}^p |F_{ij}|^2 x_j^2  = \sum_{j=1}^p \frac{1}{p} x_j^2 = \frac{\|x\|^2}{p} = \frac{1}{p}.$$

We can now use Theorem~\ref{thm:hoeffding1} with $t = \sqrt{2 \log(4p/\delta)/p}$ and obtain
$$\P \left( |y_i| >  \sqrt{ \frac{2\log(4p/\delta)}{p}} \right) \le 2\exp\left(-\displaystyle{\frac{2 \log(4p/\delta)/p}{2/p}}\right) =   \frac{\delta}{2p}.$$



Applying the union bound finishes the proof. 
\end{proof}

\begin{lemma} \label{le:par3}
Let $0<\eps<1$ and $y\in \RR^p$ satisfy  $\|y \|_\infty^2 <  \frac{2\log (4p/\delta)}{p}$ and $\|y\|=1$, then it possible to pick $d \sim 1/\eps^2 \log(p/\delta) \log(1/\delta)$ such that
%In the setting described above (and in particular assuming $\|x\|=1$), conditioned on the event that $\|y \|_\infty^2 \gtrsim  \frac{2\log (4p/\delta)}{p}$, it holds that 
$$\P\big( \|Sy\|^2 - 1 | \le \eps \big) \le  1- \frac{\delta}{2}.$$
\end{lemma}

\begin{proof}


The idea is to use Bernstein's Inequality (Theorem~\ref{thm:Bernstein}).
We write $S_{ji} = \sqrt{p/d} \delta_{ji}$, where $\delta_{ji} \in \{0,1\}$ is our random sample of the columns for row $j$. Hence for all $j$,
$\sum_{i=1}^p \delta_{ji} = 1$. We write $z:= Sy$ and compute
\begin{align*}
q_j & := z_j^2 \\
      & = \frac{p}{d} \left( \sum_{i=1}^p \delta_{ji} y_i \right)^2 \\
      & = \frac{p}{d} \left( \sum_{i} \delta_{ji} y_i^2 +   \sum_{i \neq \ell} \delta_{ji} \delta_{j\ell} y_i y_{\ell} \right) \\
       & = \frac{p}{d} \sum_{i} \delta_{ji}y_i^2.
\end{align*}
We want to bound the random variable $\|z\|^2 -1= \sum_j \left(q_j-\frac1d\right)$.  Since the $q_j$'s are independent, and $\EE\left[q_j-\frac1d\right]=0$, we can apply Berstein's Inequality (Theorem~\ref{thm:Bernstein}, provided we bound $\sigma^2$ and $a$. Firstly,
$$a \le \frac{p}{d}\|y\|_\infty^2 -\frac1d < \frac{p}{d}\|y\|_\infty^2 \leq  \frac{2\log (4p/\delta)}{d}.$$% \lesssim \frac{\log(p/\delta)}{d}.$$
and secondly,
\begin{align*}
d\sigma^2 &  \le d \E[ q_1^2]  - 1 \le d \E[ q_1^2]  = d \E \left[ \frac{p^2}{d^2} \sum_i \delta_{ji} y_i^4 \right] \\
& = \frac{p^2}{d} \frac1p\sum_i y_i^4 \\
& \leq \|y\|_\infty^2 \frac{p}{d} \sum_i y_i^2 \\
& = \frac{p}{d} \|y\|_\infty^2  \|y\|^2 \\
& = \frac{p}{d} \|y\|_\infty^2 \\
& \leq  \frac{2\log (4p/\delta)}{d}.
\end{align*}
Plugging these terms into Bernstein's Inequality (Theorem~\ref{thm:Bernstein}) gives
\begin{align*} \P \big( | \|Sy\|^2 - 1 | > \eps \big)  & = \P \Big( \Big| \sum_{j=1}^d q_j - 1 \Big| > \eps \Big) \\
& \leq 2 \exp\left(  \frac{-\eps^2}{2d\left(\frac{2\log (4p/\delta)}{d^2}\right) + \frac23\left(\frac{2\log (4p/\delta)}{d}\right)\eps  } \right)\\
& \leq 2 \exp\left(  \frac{-\eps^2}{\frac83\left(\frac{2\log (4p/\delta)}{d}\right)  } \right),\\
& = 2 \exp\left(  \frac{-3\eps^2d}{16\log (4p/\delta)  } \right),
\end{align*}
where the last inequality uses $\eps<1$.
% \lesssim \max \Big\{ e^{- \frac{c \eps^2 d}{\log(p/\delta)}}, e^{-\frac{c \eps d}{\log(p/\delta)}}\Big\}.$$
Picking $d = \frac{16}{3\eps^2} \log(4p/\delta) \log(4/\delta)$ gives the desired $\delta/2$-bound.

\end{proof}


Combining Lemmas~\ref{le:par2} and~\ref{le:par3}, union bounding over the two events, and recalling that if $\|x\|=1$ then $\|y\|=1$ since both $D$ and $F$ are unitary, establishes Theorem~\ref{th:fjl}.

Besides the potential speedup, another advantage of the FJLT  is that it requires significantly less memory  compared to storing an unstructured random projection matrix as is the case for the standard Johnson-Lindenstrauss approach. 






There are many applications where one wants to perform dimension reduction while preserving the geometry of certain sets with infinite points, such as subspaces. Via an $\epsilon$-net argument it is possible to show that random matrices that are good as Johnson-Lindenstrauss maps for $N$ points, also approximately preserve the geometry of a ($\log N$)-dimensional subspace. These go by the name of Oblivious Subspace Embeddings (see, e.g., Section 5.2.2. in~\cite{KireevaTropp} where the fast JL construction we introduced above is shown to be an Oblivious Subspace Embedding). A similar idea will be explored below (Section~\ref{s:RSVD}) where a random projection will allow us to still approximate the subspace spanned by the singular vectors corresponding to the largest singular values, resulting in a randomized SVD algorithm that is workhorse of large scale linear algebra. In Chapter~\ref{c:probability-gaussiananalysis} we will develop analogues of the Johnson-Lindenstrauss Lemma for more general sets,  through the celebrated Gordon's espace through a mesh theorem.



\section{Randomized SVD}\label{s:RSVD}

Randomized linear algebra has gained significant relevance in recent years due to its ability to efficiently handle large-scale data problems that traditional methods struggle with.  
The Johnson-Lindenstrauss dimension reduction discussed in the previous section is one prominent example. The techniques encountered in the Johnson-Lindenstrauss transform naturally lend themselves as key ingredient of the \emph{randomized SVD},  a method to compute an approximation of the SVD of a matrix. It is especially useful when the matrix is large and dense and can be well approximated by a low-rank matrix. 

In this section we will give a detailed description of the randomized SVD.
It is based on using random projections to reduce the dimensionality of the matrix before performing a more traditional SVD on the reduced matrix.


Recall from Chapter~\ref{c:svd} that  the SVD of  a matrix $A\in\R^{m\times n}$ is given by
\begin{equation}\label{eq:SVDrand}
 A = U \Sigma V^{\ast},
\end{equation}
where $U \in \mathbb{R}^{m \times m}$ is a unitary matrix containing the left singular vectors,
$\Sigma \in \R^{m \times n}$ is a diagonal matrix with singular values (in decreasing order) on the diagonal,
and $V \in \mathbb{R}^{n \times n}$ is a unitary matrix  containing the right singular vectors.

In practice $A$ is often large, and we are interested in only a small number $s$ of the largest singular values and corresponding singular vectors, giving a rank-$k$ approximation of $ A $:
$$
A \approx U_k \Sigma_k V_k^{\ast}
$$
where $ U_k \in \mathbb{R}^{m \times k} $, $ \Sigma_k \in \R^{k \times k} $, and $ V_k \in \mathbb{R}^{n \times k} $.


We define the projection $P_k$ onto the the span of the leading $k$ left singular vectors of $A$ via
$$
P_k = \sum_{i=1}^k u_i u_i^{\ast}.
$$
We know from Theorem~\ref{th:rank_k_Frob} that the best rank-$k$ approximation of $A$ with respect to the Frobenius norm is given by
%$A_s = P_sA $ and the resulting error is
$$\min_{\operatorname{rank}(B) \leq k} \|A - B \|_F = \|A - P_k A \|_F =  \sqrt{\sum_{i > k} \sigma_i^2}.$$
Thus, we can compute the best approximation by projecting the target matrix onto the $k$-dimensional
subspace that captures ``most of the action'' of the matrix $A$.

In many applications it suffices to compute the best approximation with moderate accuracy. Hence, following~\cite{HMT11}, we let $s=k+p$ for small $p$, and seek a rank-$s$ approximation $\hat{A}_s$ that is comparable to the best rank-$k$ approximation, i.e.,
$$\|A - \hat{A}_s\|_F^2  \le (1+\eps) \|A - P_k A \|_F^2 = (1+\eps)\sum_{i > k} \sigma_i^2 $$
for some chosen tolerance $\eps > 0$.

The aim of the Randomized SVD (RSVD) is essentially to generate an efficient approximation for the best rank-$k$ approximation of the matrix $P_k A$. 
We  proceed in two steps. In the first step we compute a low-dimensional subspace that captures the range of $A$. In the second step we project $A$ onto the subspace and compute the (standard SVD) of the resulting smaller matrix.
Regarding the first step, recall that $P_k$ is the orthogonal projection  onto the $k$
leading left singular vectors of $A$. 
But since we do not know the subspace spanned by $P_k$, a sensible choice, as suggested by the  Johnson-Lindenstrauss Lemma,  is to draw some {\em probing} random vectors from a rotationally invariant distribution to estimate the directions of the vectors $u_i$ that make up $P_k$. We collect these probing random vectors in a matrix $\Psi$.


Hence, let $ \Psi \in \R^{n \times (k + p)} $ be a  random matrix, where $ k $ is the desired number of singular values/vectors,
and $ p $ is an oversampling parameter (typically $ p $ is small, like 5 or 10).  A promising choice for $ \Psi $ is to select
each entry of $ \Psi $ as a sample from a standard normal distribution $ \mathcal{N}(0, 1) $.
We then form the matrix $ Y \in \R^{m \times (k + p)} $ as
$$
Y = A \Psi.
$$
This matrix $ Y $ captures the action of $ A $ on a random subspace of dimension $ k + p $. We will see (in Theorem~\ref{th:randSVD}) that, with high probability, the subspace spanned by the columns of $ Y $ contains a good approximation of the subspace spanned by the top $ k $ left singular vectors of $ A $.\footnote{In many applications the matrix $A$ is implicit and one has access to it via matrix-vector products, if this is the case computing $Y$ can be done with $k+p$ matrix-vector products.}
Next, we find an orthonormal basis for the range of $ Y $. Using the standard QR decomposition, we compute
$$
Y = Q R,
$$
where $ Q \in \mathbb{R}^{m \times (k + p)} $ has orthonormal columns, and $ R \in \R^{(k + p) \times (k + p)} $ is an upper triangular matrix.
Now, we project $ A $ onto the lower-dimensional subspace spanned by the columns of $ Q $ by computing
$$
B = Q^{\ast} A.
$$
Here, $ B \in \R^{(k + p) \times n} $ is a much smaller matrix that approximates $ A $ (in the sense that $QB$ approximates $A$).

Next, we compute the exact SVD of the smaller matrix $ B $
$$
B = U_B \Sigma_B V_B^{\ast}.
$$
where $U_B \in \mathbb{R}^{(k + p) \times (k + p)} $, $ \Sigma_B \in \R^{(k + p) \times n} $, and $ V_B \in \R^{n \times  (k + p)} $.
Finally, the approximate SVD of $ A $ is given by
$$
A \approx \tilde{ U} \tilde{ \Sigma} \tilde{V}^{\ast}
$$
where $ \tilde{U} = Q U_B \in \mathbb{R}^{m \times (k + p)} $, $ \tilde{\Sigma} = \Sigma_B \in \R^{ (k + p) \times  (k + p)} $, and $ \tilde{V} = V_B \in \R^{n \times  (k + p)} $.



We summarize these steps in Algorithm~\ref{alg:RandSVD}.

\begin{algorithm}
%\SetAlgoLined
\caption{Randomized SVD}
\label{alg:RandSVD}
%\begin{algorithmic}[1]
%\begin{flushleft}
\textbf{Input:} Matrix  $A \in\R^{m\times n}$, target rank $k>0$, oversampling factor $p>0$; set $s=k+p$.
\begin{enumerate}
\item Generate a random matrix $\Psi \in \mathbb{R}^{N \times (k+p)}$ 
\item  Compute the product $Y = A \Psi$ 
\item Compute the economy-sized QR factorization $Y = Q R$, with $Q \in\R^{m\times (k+p)}$ and  $R \in\R^{(k+p)\times (k+p)}$ 
\item Compute the product  $B = Q^{\ast} A$ 
\item Compute the economy-sized SVD  $B =  U_B \Sigma_B V_B^{\ast}$, with $U_B \in\R^{m\times (k+p)}, \Sigma_B \in\R^{(k+p)\times (k+p)}$ and  $V_B \in\R^{n \times (k+p)}$ 
\item Set $\tilde{U} = Q U_B, \tilde{\Sigma} = \Sigma_B, \tilde{V} = V_B.$ 
\end{enumerate}
\textbf{Output:} Orthogonal $\tilde{U} \in \R^{m \times (k+p)}$, orthogonal $\tilde{V} \in \R^{n \times (k+p)}$, and diagonal $\tilde{\Sigma} \in  \R^{(k+p) \times (k+p)}$ 
such that $A \sim \hat{A}_s:= \tilde{U} \tilde{\Sigma} \tilde{V}^{\ast}$ 

%\end{flushleft}
%\end{algorithmic}
\end{algorithm}


The random matrix $ \Psi $ used in the random projection step ensures that with high probability, the subspace spanned by the columns of $ Y = A \Psi $ contains a good approximation to the dominant singular subspace of $ A $. 



\begin{theorem}[\cite{HMT11}]\label{th:randSVD}
Consider a matrix $A \in \R^{m \times n}$ with singular values $\sigma_1 \ge \sigma_2 \ge \sigma_3 \ge \cdots$. Fix the target rank  $k \le \min(m,n)$. When $s :=k+p \ge k+2$, the randomized SVD method  produces a random rank-$s$
approximation $\hat{A}_s$ that satisfies
\begin{equation}
\label{randSVDbound}
\E[ \|A - \hat{A}_s \|_F ] \le \left(1 + \frac{k}{p-1}\right)^\frac{1}{2} \Big( \sum_{j > k}\sigma_{j}^2\Big)^\frac{1}{2}.
\end{equation}
\end{theorem}

The error on the right hand side consists of the best rank-$k$ approximation error and an additional factor due to using a randomized projection instead of the optimal projection $P_k$.

For the proof of Theorem~\ref{th:randSVD} we need some preparation. In what follows it will be convenient to rewrite the SVD of $A$, cf.~\eqref{eq:SVDrand}, in block matrix form:
\begin{equation}
A = \begin{bmatrix}  U_k & U_\perp \end{bmatrix} \begin{bmatrix}  \Sigma_k & 0 \\ 0 & \Sigma_\perp \end{bmatrix} \begin{bmatrix}  V_k^{\ast} \\ V_\perp^{\ast} \end{bmatrix},
\end{equation}
where $ U_\perp \in \R^{m \times m-k}$ denotes the orthogonal complement of $U_k \in \R^{m \times k}$ with respect to $U \in \R^{m \times m}$.
Furthermore, we set 
$$\Psi_k : = V_k^\ast \Psi \qquad \text{and} \qquad \Psi_\perp : = V_\perp^\ast \Psi .$$
Here, $\Psi_k$ captures the alignment of $\Psi$ with the leading right singular vectors of $A$, while  $\Psi_\perp$ captures
the alignment with the trailing right singular vectors of $A$. 
The matrix $P_Y$ denotes the orthogonal projection onto the range of $Y$. Since this projection is unique, we have that $P_Y = Q Q^{\ast}$.


The following deterministic error bound (the bound does not
require $\Psi$ to be Gaussian) will be instrumental:

\begin{lemma}\label{le:SVDerrorbound}
We use the same notation as in Theorem~\ref{th:randSVD}. Let $\Psi \in \R^{n \times (k+p)}$ be an arbitrary matrix with 
 $\rank(\Psi)=k$. Then the rank-$s$ approximation $\hat{A}_s$ computed by Algorithm~\ref{alg:RandSVD} satisfies
\begin{equation}\label{RSVDdeterministic}
\| A - \hat{A}_s\|^2_F \le \|\Sigma_\perp\|_F^2 + \|\Sigma_k \Psi_k \Psi_\perp^{\dagger}\|_F^2.
\end{equation}
\end{lemma}


\begin{proof}
We first fix the target rank $k \le \min{\{m,n\}}$ and the oversampling factor $p$ and set $s=k+p$. 
%It will be convenient to isolate the $k$ leading singular directions in the SVD of $A$. 
Denoting $\tilde{A} = U^{\ast} A$ and $\tilde{Y}= \tilde{A}\Psi$, we obtain
\begin{equation}\label{RSVDdeterministic1}
\| A - \hat{A}_s\|^2_F = \| (I - P_Y)A \|^2_F = \| U^{\ast} (I - P_Y) U\tilde{A} \|^2_F = \| ( I - P_{\tilde{Y}} ) \tilde{A} \|^2_F.
\end{equation}
Here, the first equality uses the fact that the Frobenius norm is unitarily invariant, whereas in the second equality we have substituted $P_{\tilde{Y}} = U^{\ast} P_Y U$.
Next, we set $Z=\tilde{Y} \Psi_k^{\dagger} \Sigma_k^{-1} =: [I \,\,\, F]^{\ast}$ where $F = \Sigma_\perp \Psi_\perp \Psi_k^{\dagger} \Sigma_k^{-1}$.
Note that $\range(Z) \subseteq \range(\tilde{Y})$. This in turn yields
\begin{equation}\label{RSVDdeterministic2}
\| (I-P_{\tilde{Y}} )\tilde{A}\|_F^2 \le \| (I - P_Z)\tilde{A} \|_F^2 = \tr \big( \tilde{A}^{\ast} (I - P_Z)\tilde{A} \big) = \tr \big(\Sigma (I - P_Z) \Sigma \big).
\end{equation}
Note that $P_Z$ has the expansion
$$
P_Z =  \begin{bmatrix}  I \\ F  \end{bmatrix}  \begin{bmatrix}  I + F^{\ast} F  \end{bmatrix}^{-1}  \begin{bmatrix}  I \\ F  \end{bmatrix}^{\ast}.
$$
Since $P_Z$ is a projection, we have that $I - P_Z \succ 0$. After some linear algebra calculations (left to the reader), we obtain
$$ 0 \prec I - P_X \preceq \begin{bmatrix}  F^{\ast}F & B \\ B^{\ast} &  I \end{bmatrix},$$
where $B = -(I+F^{\ast}F)^{-1} F^{\ast}$. Multiplying both sides by $\Sigma$ gives
\begin{equation}\label{RSVDdeterministic4}
\Sigma(I-P_Z)\Sigma  \preceq \begin{bmatrix}  \Sigma_k F^{\ast}F \Sigma_k & \Sigma_k B \Sigma_\perp \\ \Sigma_\perp B^{\ast} \Sigma_k & \Sigma_\perp \Sigma_\perp \end{bmatrix}.
\end{equation}
Combining~\eqref{RSVDdeterministic1},  \eqref{RSVDdeterministic2} and ~\eqref{RSVDdeterministic4} gives
\begin{flalign*}
\| A- \hat{A}_s \|_F^2 & = \|(I-P_{\tilde{Y}}) \tilde{A} \|_F^2  \\
                                  & \le \tr \big( \Sigma (I-P_Z) \Sigma \big) \\
                                  & \le \tr \big(\Sigma_k F^{\ast} F \Sigma_\perp \big) + \tr \big( \Sigma_\perp \big) \\
                                  & \le \|\Sigma_\perp \Psi_\perp \Psi_k^{\dagger} \|_F^2 + \|\Sigma_\perp\|_F^2,
\end{flalign*}
where in the penultimate inequality we have used the fact that the trace is monotone on the cone of positive semidefinite matrices.

\end{proof}



Lemma~\ref{le:SVDerrorbound} holds for {\em any} rank-$k$ matrix $\Psi$.  As pointed out in~\cite{KireevaTropp}, the first term on
the right-hand side of~\eqref{RSVDdeterministic}  captures the trailing singular values of the matrix, which is independent of the
specific choice of $\Psi$, since this loss occurs for every rank-$k$ approximation. The second term is the error caused by the
dimension reducing matrix $\Psi$.
We want the component $\Psi_k$ in the leading directions to be well conditioned (ideally, an
identity matrix). We want the component $\Psi_\perp$ in the trailing directions to be as small as possible.


The bound~\eqref{RSVDdeterministic} suggests that
in order to establish Theorem~\ref{th:randSVD} we need to get a handle on $ \Psi_k \Psi_\perp^{\dagger}$ when $\Psi_{ij}
\overset{\text{i.i.d.}}{\sim} {\mathcal N}(0,1).$ 
To that end, the following lemma will be helpful.



\begin{lemma}\label{pr:Gaussnorm2}
Let  $G \in \R^{m \times n}$ be a standard Gaussian random matrix. 
 Assume $m \ge 2$ and $n-m \ge 2$. Then
\begin{equation}\label{eq:gaussnorm3}
\E [ \| G^{\dagger} \|_F^2 ] = \frac{m}{n-m-1}.
\end{equation}
\end{lemma}

\begin{proof}
We follow the proof of the corresponding result in the appendix of~\cite{HMT11}.
 First, note that $(G G^{\ast})$ is invertible almost surely. We compute
$$  \| G^{\dagger} \|_F^2 = \trace [(G^{\dagger})^{\ast} G^{\dagger}] = \trace [(G G^{\ast})^{-1}].$$
The random matrix $(G G^{\ast})^{-1}$ follows the inverted Wishart distribution. Hence, we can use formula (12) on page 97 of~\cite{muirhead2009aspects} for the expected trace and obtain~\eqref{eq:gaussnorm3}.



\end{proof}

With these lemmata in place, we are now ready to prove Theorem~\ref{th:randSVD}.
\begin{proof}[of Theorem~\ref{th:randSVD}]
Since $\Psi$ is chosen to be a standard Gaussian matrix and the Gaussian distribution is invariant under unitary transformations, it follows that $\Psi_k=V^{\ast}_k \Psi$ and $\Psi_\perp =V^{\ast}_\perp \Psi$ are also standard Gaussian. 
%Also,  $\Psi_k$ and $\Psi_\perp$ are statistically independent.
Moreover, note that matrix $\Psi_\perp$ has full rank with probability one, since the rows of a (fat) Gaussian matrix are almost surely in general position.

We consider the bound established in Lemma~\ref{le:SVDerrorbound}, apply expectation to~\eqref{RSVDdeterministic} and use H\"{o}lder's inequality to obtain
\begin{equation}\label{RSVDbound4}
\E [ \| A - \hat{A}_s\|_F]  
\le \big(\E [ \| (I-P_Y)A\|_F^2] \big)^{1/2}
\le \Big( \|\Sigma_\perp\|_F^2 + \E [\|\Sigma_k \Psi_k \Psi_\perp^{\dagger}\|_F^2 ]  \Big)^{1/2}.
\end{equation}
Using the law of total expectation as well as the independence of
$\Psi_k$ and $\Psi_\perp$ we estimate
\begin{align}%\label{RSVDbound6}
\E [ \|\Sigma_k \Psi_k \Psi_\perp^{\dagger}\|_F^2 ]  & = \E_{\Psi_k} \big[ \E [ \| \Sigma_\perp \Psi_\perp \Psi_k^{\dagger} \|_F^2 | \Psi_k ] \big] \notag \\ 
& \le  \E \big [ \| \Sigma_\perp \|_F^2 \,  \|\Psi_k^{\dagger} \|_F^2 \big]  \notag \\
& \le  \| \Sigma_\perp \|_F^2 \, \E \big[ \|\Psi_k^{\dagger} \|_F^2 \big]  \notag \\
& \le  \frac{k}{p-1}  \sum_{j > k} \sigma_j^2, 
\label{RSVDbound6}
\end{align}
where the second  inequality follows from Lemma~\ref{pr:Gaussnorm2}.

Combining now~\eqref{RSVDbound4} with~\eqref{RSVDbound6} establishes the bound~\eqref{randSVDbound}.

\end{proof}


\subsection{Computational complexity of the Randomized SVD}

What is the computational complexity of the RSVD? The dominant costs for the RSVD are:
\begin{enumerate}
\item ${\mathcal O}(m n (k + p)) $ for the matrix-matrix multiplication $ A \Psi $, assuming that $\Psi$ is an unstructured, non-sparse matrix (like a Gaussian matrix).
\item $ {\mathcal O}(m (k + p)^2) $ for the QR decomposition of $ Y $.
\item ${\mathcal O}(mn (k + p))$  for computing $Q^{\ast} A$.
\item  ${\mathcal O}((k + p)^2 n) $ for computing the  SVD of the small matrix $ B $.
\end{enumerate}
Thus, the overall complexity of the randomized SVD is ${\OOO}(m n (k + p)) $, which is significantly smaller than
the $ {\OOO}(m n^2) $ cost of the classical SVD when $ k $ is much smaller than $ n $.

We can further reduce the cost of the matrix-matrix product  $ A \Psi $ by introducing structure or sparsity into $\Psi$. We already encountered one such choice in Definition~\ref{FastJL} in the form of the Fast Johnson-Lindenstrauss transform. The improved computational efficiency usually comes at the cost of additional log-factors in terms of required samples to guarantee the same accuracy.


When the singular values of $A$ decay slowly, Algorithm~\ref{alg:RandSVD} is not very accurate in the regime of interest, namely when $k \ll n$. The reason is that to achieve a small approximation error the matrix $Q$ must align well with the first $k + p$ left singular vectors of $A$. As a consequence, the product $A\Psi$ must reveal
the singular vectors, which in turn would require the singular vectors associated with small singular values not to interfere with the calculation. 

How can we remedy this issue? We can take inspiration from the power iteration. 
By raising our matrix to a power $q$ via $(AA^{\ast} )^q A$ prior to applying it to $\Psi$ we can enhance the decay of the 
spectrum of our matrix without effecting the direction of the eigenvectors. Hence, we replace in Algorithm~\ref{alg:RandSVD} the step $Y = A \Psi$ by $Y = (AA^{\ast} )^q A \Psi$.
A small power of $q$ (such as $q=1$ or 2) is often sufficient. In this modified RSVD  it is also advisable to increase the oversampling factor $p$. We refer to~\cite[Corollary 10.10]{HMT11} for an approximation error bound for this modified RSVD as well as to~\cite{tropp2023randomized} for a more detailed discussion.

The key ideas behind the RSVD also form the main ingredients for some other algorithms, such as randomized least squares and randomized Choleksy factorization, see~\cite{rokhlin2008fast,HMT11,KireevaTropp}.

\section{Matrix sketching}\label{s:sketching}

The methods we encountered in this chapter are closely intertwined with the concept of {\em matrix sketching}~\cite{liberty2013simple,mahoney2011randomized,rudelson2007sampling,HMT11}, a family of randomized techniques for constructing a low-dimensional representation of a large matrix while approximately preserving its essential linear-algebraic structure. 

Given a data matrix $A \in \mathbb{R}^{n \times d}$, the goal is to replace $A$ by a much smaller matrix—called a {\em sketch}---that enables accurate approximation of quantities such as norms, inner products, or dominant singular subspaces. 

Two complementary formulations are common. In the {\em projection-based} viewpoint, one forms a sketch $SA$, where $S \in \mathbb{R}^{m \times n}$ is a random embedding (e.g., Gaussian or subsampled randomized Hadamard), generalizing the Johnson–Lindenstrauss lemma from vectors to entire subspaces~\cite{woodruff2014sketching,tropp2017practical}. 

In the {\em sampling-based} viewpoint, one constructs a sketch by randomly selecting and rescaling columns (or rows) of $A$, often using leverage-score sampling to retain the most informative directions~\cite{frieze2004fast,mahoney2009cur,drineas2006subspace}. Both approaches can be understood as approximately preserving the geometry of the column space of $A$. 

Matrix sketching does not only underlie the randomized SVD, it provides a principled way to trade controlled approximation error for potentially substantial reductions in computational and memory costs. As such, it enables distributed linear algebra~\cite{woodruff2014sketching} and it is the fundamental mathematical too that makes {\em streaming} possible for high-dimensional data~\cite{harvey2008sketching,mahoney2011randomized}. Here, streaming refers to a computational model in which data are processed sequentially, in one (or very few) passes, using limited memory, without storing the entire dataset~\cite{muthukrishnan2005data,woodruff2014sketching}. Streaming is essential when
the dataset is too large to store in memory,
data arrive continuously (logs, sensor data, social networks), or computation must be done online.

Let us explore some exemplary results in sketching beyond the Johnson-Lindenstrauss transform and the randomized SVD.

\begin{definition}
A $(1\pm \eps)$ $\ell_2$-subspace embedding for the column space of an $n \times d$ matrix $A$ is a matrix $S$ for which for all $x \in \R^d$
$$
(1-\varepsilon)\|Ax\|^2 \le \|SAx\|^2 \le (1+\varepsilon)\|Ax\|^2
\quad \text{for all } v \in \mathbb{R}^d.
$$
\end{definition}

We will occasionally abuse notation and say that $S$ is an $\ell_2$-subspace embedding for $A$ itself, although it should be clear from the definition that this property is independent from the choice of a particular basis for the representation
of the column space of $A$.


\begin{theorem}[$\ell_2$-subspace embedding]\label{th:subspaceembedding} 
Let $U \subset \mathbb{R}^n$ be a fixed subspace of dimension $r$.
Let $S \in \mathbb{R}^{m \times n}$ be a random matrix such that for any fixed vector $x \in \mathbb{R}^n$,
$$
\mathbb{P}\left( \left| \|Sx\|^2 - \|x\|^2 \right| > \varepsilon \|x\|^2 \right)
\le 2e^{-c\varepsilon^2 m}
$$
If
$$
m \ge C \frac{r + \log(1/\delta)}{\varepsilon^2},
$$
then with probability at least $1-\delta$,
$$
(1-\varepsilon)\|x\|^2 \le \|Sx|^2 \le (1+\varepsilon)\|x\|^2
\quad \forall x \in U.
$$
\end{theorem}

\begin{proof}
By homogeneity, it suffices to prove the inequality for all
$$
x \in \mathbb{S}_U := \{x \in U : \|x\| = 1\}.
$$
There exists an $\eta$-net $\mathcal{N} \subset \mathbb{S}_U$ with
$
|\mathcal{N}| \le \left(1+\frac{2}{\eta}\right)^r.
$
Choose $\eta = \frac{\varepsilon}{4}$, then
$
|\mathcal{N}| \le \left(\frac{8}{\varepsilon}\right)^r
 \le
e^{C r}.
$
For each fixed $x \in \mathcal{N}$,
$$
\P \left(
\bigl| \|Sx\|^2 - 1 \bigr| > \frac{\varepsilon}{2}
\right)
 \le
2 e^{-c \varepsilon^2 m}.
$$
By the union bound,
$$
\P \left(
\exists x \in \mathcal{N} :
\bigl| \|Sx\|^2 - 1 \bigr| > \frac{\varepsilon}{2}
\right)
\le
|\mathcal{N}| \cdot 2 e^{-c \varepsilon^2 m}.
$$
Using the bound on the size of the net, $|\mathcal{N}| \le e^{C r}$,
we obtain
$$
\P \left(\exists x \in \mathcal{N} :
\bigl| \|Sx\|^2 - 1 \bigr| > \frac{\varepsilon}{2} \right)
\le 2 e^{ C r - c \varepsilon^2 m }.
$$
If
$$
m \ge C \frac{r + \log(1/\delta)}{\varepsilon^2},
$$
then
$$
Cr - c \varepsilon^2 m \le -\log(2/\delta).
$$

Hence
$$
\P \left(
\forall x \in \mathcal{N} :
\bigl| \|Sx\|^2 - 1 \bigr| \le \frac{\varepsilon}{2}
\right)
\ge 1-\delta.
$$
From now on, we condition on this event.

Let $y \in \mathbb{S}_U$ and choose $x \in \mathcal{N}$ such that
$$
\|y - x\| \le \eta = \frac{\varepsilon}{4}.
$$
We write
$$
\|Sy\|
\le
\|Sx\| + \|S(y-x)\|.
$$

Using concentration again for the fixed vector $y-x$,
$$
\|S(y-x)\|^2
\le
(1+\varepsilon)\|y-x\|^2
\le
(1+\varepsilon)\frac{\varepsilon^2}{16}.
$$

Combining with $\|Sx\|^2 \le 1+\varepsilon/2$, one obtains
$\|Sy\|^2 \le 1+\varepsilon$.
A symmetric argument gives the lower bound
$\|Sy\|^2 \ge 1-\varepsilon$.

Hence, with probability at least $1-\delta$,
$$
(1-\varepsilon)\|x\|^2
\le
\|Sx\|^2
\le
(1+\varepsilon)\|x\|^2
\quad \forall x \in U.
$$
 
\end{proof}

It is easy to see from our previously derived concentration bounds that a Gaussian random matrix and a Subsampled Randomized Hadamard (or Fourier) Transform as introduced in Definition~\ref{FastJL}
satisfy the conditions of Theorem~\ref{th:subspaceembedding}.

Another interesting and useful construction of a sketching matrix is the so-called {\em CountSketch}~\cite{charikar2002finding,thorup2012tabulation}. CountSketch is extremely sparse, making it an appealing choice for processing massive datasets or high-speed data streams. It is often used to estimate frequencies of elements in data streams.

A CountSketch matrix $S \in \mathbb{R}^{m \times n}$ is defined by two independent {\em hash functions}:
\begin{enumerate}
\item $h: \{1, \dots, n\} \to \{1, \dots, m\}$: This {\em bucket hash function} maps each column index to one of the $m$ rows.
\item $g: \{1, \dots, n\} \to \{+1, -1\}$: This {\em sign hash function} assigns a random sign to each column.
\end{enumerate}

The entries of $S$ are constructed as follows:
$$S_{i,j} = \begin{cases} g(j) & \text{if } h(j) = i, \\ 0, & \text{otherwise}. \end{cases}$$

The CountSketch construction enjoys the following useful properties: (i)~Sparsity: Each column of $S$ has exactly one non-zero entry. (ii)~Efficiency: To compute $Sx$, we do not perform a standard matrix-vector product. Instead, for each entry $x_j$, we simply update one row of the sketch: $sketch[h(j)] \leftarrow sketch[h(j)] + g(j)x_j$. (iii)~Time Complexity: Computing $SA$ for $A \in \mathbb{R}^{n \times d}$ takes only $\OOO(\text{nnz}(A))$ (the number of non-zero entries in $A$) operations.

A typical realization of $S$ might look like this:
$$S = \begin{bmatrix} 0 & 0 & 0 & 0 & -1 & 0 \\ +1 & 0 & 0 & 0 & 0 & 0 \\ 0 & 0 & 0 & +1 & 0 & 0 \\ 0 & -1 & 0 & 0 & 0 & +1 \\ 0 & 0 & -1 & 0 & 0 & 0 \\ 0 & 0 & 0 & 0 & 0 & 0 \end{bmatrix}.$$
In this specific example, the two underlying hash functions performed the following actions:
(i)~The bucket hash ($h$):
$h(1)=2, h(2)=4, h(3)=5, h(4)=3, h(5)=1, h(6)=4$.
Note that row 4 received two assignments (from columns 2 and 6, respectively), while row 6 is empty. This ``collision'' is expected in hashing. (ii)~The sign hash ($g$):
$g(1)=+1, g(2)=-1, g(3)=-1, g(4)=+1, g(5)=-1, g(6)=+1$.


Before proving subspace embeddings, we establish that CountSketch preserves the inner product.

\begin{lemma}\label{le:countsketchinner}
For any two vectors $x, y \in \mathbb{R}^n$, the inner product of their sketches via the CountSketch matrix $S$ is an unbiased estimator of their true inner product:
$$\E[\langle Sx, Sy \rangle] = \langle x, y \rangle.$$
\end{lemma}

The proof is left to the reader in Exercise~\ref{ex:countsketch}.


\begin{theorem}\label{thm:sketch-countsketch}
Let $U \subset \mathbb{R}^n$ be a subspace of dimension $d$. If the number of rows $m$ satisfies
$$m \ge \frac{2d^2}{\delta\epsilon^2},$$
then with probability $1-\delta$
$$(1-\epsilon)\|x\|^2 \le \|Sx\|^2 \le (1+\epsilon)\|x\|^2 \quad \forall x \in U.$$
\end{theorem}

\begin{proof}
A matrix $S$ is an $\epsilon$-subspace embedding for the column space of $U$ if and only if the eigenvalues of $(SU)^T (SU)$ are within $[1-\epsilon, 1+\epsilon]$. This is equivalent to
$$\| (SU)^T (SU) - I_d \| \le \epsilon.$$
Since the spectral norm is upper-bounded by the Frobenius norm, it suffices to show that 
$\mathbb{P}(\| (SU)^T (SU) - I_d \|_F^2 > \epsilon^2) \le \delta$.

Let $M = (SU)^T (SU) - I_d$. The entries of the product $(SU)^T (SU)$ are
$$((SU)^T (SU))_{jk} = \langle Su_j, Su_k \rangle$$
where $u_j$ and $u_k$ are the $j$-th and $k$-th columns of $U$.
From Lemma~\ref{le:countsketchinner} we know $\E[\langle Su_j, Su_k \rangle] = \langle u_j, u_k \rangle$.
Since $U$ has orthonormal columns, $\langle u_j, u_k \rangle = \delta_{jk}$. Thus, $\E[M] = 0$.

The squared Frobenius norm is the sum of the squares of the entries
$$\E[\|M\|_F^2] = \sum_{j=1}^d \sum_{k=1}^d \E[((SU)^T (SU) - I_d)_{jk}^2].$$
We expand the expectation for a single pair $(j, k)$. Let $X_{jk} = \langle Su_j, Su_k \rangle$. We seek $\Var(X_{jk})$ because $\E[(X_{jk} - \delta_{jk})^2] = \Var(X_{jk})$.
Using the definition of the CountSketch matrix, we can expand
$X_{jk}$ as follows:
$$X_{jk} = \sum_{a=1}^n \sum_{b=1}^n g(a)g(b) \mathbf{1}_{h(a)=h(b)} U_{aj} U_{bk}$$
Using the property of random signs $g$, the expectation of $X_{jk}^2$ survives only when indices match in pairs. After some algebraic simplification using the independence of $h$ and $g$ we obtain
$$\Var(\langle Su_j, Su_k \rangle) = \frac{1}{m} \left( \|u_j\|^2 \|u_k\|^2 + \langle u_j, u_k \rangle^2 \right) - \frac{2}{m} \sum_{a=1}^n u_{aj}^2 u_{ak}^2$$
Summing over all $j, k \in \{1, \dots, d\}$ gives
$$\E[\|M\|_F^2] = \sum_{j,k} \Var(X_{jk}) \le \frac{1}{m} \sum_{j,k} (\|u_j\|^2 \|u_k\|^2 + \langle u_j, u_k \rangle^2)$$
Since $\|u_j\|^2 = 1$ and $\sum_{j,k} \langle u_j, u_k \rangle^2 = \|U^T U\|_F^2 = \|I_d\|_F^2 = d$, we get
$$\E[\|M\|_F^2] \le \frac{d^2 + d}{m} \le \frac{2d^2}{m}.$$

Applying Chebyshev's inequality gives
$$\mathbb{P}(\|M\|_F^2 > \epsilon^2) \le \frac{\E[\|M\|_F^2]}{\epsilon^2} \le \frac{2d^2}{m \epsilon^2}$$
To ensure this probability is at most $\delta$, we set
$$\frac{2d^2}{m \epsilon^2} = \delta \implies m = \frac{2d^2}{\delta \epsilon^2}$$
This completes the proof. 
\end{proof}


The attentive reader will notice a couple of places where this argument is rather suboptimal: the tail bound is obtained with Chebyshev's inequality, and the quantity controlled is the Frobenius norm rather than the spectral norm (indeed, many $d\times d$ random matrices tend to have Frobeniues norm that is about $\sim\sqrt{d}$ times larger than the spectral norm). Using random matrix tools similar to the ones we will develop in Chapter~\ref{c:probability-matrixconcentration}, it is possible to significantly improve Theorem~\ref{thm:sketch-countsketch} (see e.g.~\cite{nelson2013osnap}). 
In any case, while CountSketch is efficient, it requires a larger sketch size than Gaussian matrices to achieve the same $\ell_2$-subspace embedding guarantee. But CountSketch is the preferred choice when $n$ and $d$ are so large that even a single dense matrix multiplication is prohibitive.










    

\bigskip






The following theorem illustrates the use of sketching to speed up the solution of large-scale least squares problems, as they arise for example in high-dimensional linear regression.


\begin{theorem}[Approximate regression via sketching]
Let $A \in \mathbb{R}^{n\times d}$, $y \in \mathbb{R}^n$, and let $S$ be an $\ell_2$-subspace embedding for $\mathrm{Range}([A; y])$ (where $\mathrm{Range}([A; y])$ denotes the subspace spanned by the columns of $A$ and the vector $y$).
Let
$$
\tilde{z} = \underset{z}{\argmin} \, \|SAx - Sy\|.
$$
Then
$$
\|A\tilde{x} - y\| \le (1+\varepsilon)\min_x \|Ax - y\|.
$$
\end{theorem}

\begin{proof}
Let
$$
r(x) := Ax - y.
$$
The least-squares solution satisfies
$$
x^* = \arg\min_x \|r(x)\|.
$$

Since $r(x) \in \range ([A; y])$, the $\ell_2$-subspace embedding implies
$$
(1-\varepsilon)\|r(x)\|
\le \|Sr(x)\|
\le (1+\varepsilon)\|r(x)\|
\quad \forall x.
$$

By definition of $\tilde{x}$,
$$
\|S r(\tilde{x})\| \le \|S r(x^*)\|.
$$

Applying the embedding bounds yields
$$
(1-\varepsilon)\|r(\tilde{x})\|
\le (1+\varepsilon)\|r(x^*)\|.
$$

Dividing by $1-\varepsilon$,
$$
\|A\tilde{x} - y\|
\le \frac{1+\varepsilon}{1-\varepsilon} \|Ax^* - y\|.
$$

For $0 < \varepsilon \le 1/3$,
$$
\frac{1+\varepsilon}{1-\varepsilon} \le 1 + 3\varepsilon,
$$
which yields the desired bound (up to rescaling $\varepsilon$).
     
\end{proof}



\subsection{Fast approximate matrix multiplication via sketching?}

Computing the product of two matrices is one of the most common and basic procedures in numerical linear algebra.
Let $A,\,B$ be two arbitrary $n\times n$ matrices. The cost for the classical way (the standard inner product method) to compute the product $AB$  is ${\mathcal O}(n^3)$.

In 1969 Strassen took the linear algebra community by surprise when he showed that the standard matrix
multiplication algorithm is not optimal from a computational viewpoint by presenting an algorithm that uses  only
${\OOO}(n^{2.808})$ operations~\cite{Strassen1969}.
Strassen's algorithm proceeds as follows.
Let us decompose the product of $A$ and $B$ into blocks
in the following way
$$
\begin{bmatrix} A_{11}&A_{12} \\ A_{21}&A_{22} \end{bmatrix}
\begin{bmatrix} B_{11}&B_{12} \\ B_{21}&B_{22} \end{bmatrix} =
\begin{bmatrix}  C_{11}&C_{12} \\ C_{21}&C_{22} \end{bmatrix},
$$
where  each block is of size $\frac{n}{2}\times\frac{n}{2}$. Following the standard
way of multiplication, this requires eight $\left(\frac{n}{2}\times\frac{n}{2}\right)$ matrix multiplications
and four additions of $\frac{n}{2}\times\frac{n}{2}$ matrices.
Strassen devised a formula that computes this block matrix multiplication using only {\em seven}
$\frac{n}{2}\times\frac{n}{2}$ matrix multiplication and eighteen $\frac{n}{2}\times\frac{n}{2}$ additions.
The cost for multiplying two $\frac{n}{2}\times\frac{n}{2}$ matrices is ${\mathcal O}\left(\frac{n}{2}\right)^3$ and the cost
for adding two $\frac{n}{2}\times\frac{n}{2}$ matrices is ${\mathcal O}\left(\frac{n}{2}\right)^2$. Since matrix addition is cheaper than matrix multiplication, there is a net gain in computational cost.
Proceeding recursively, with blocks of size $\frac{n}{2}$, $\frac{n}{4}$, \ldots, (up to blocks of size 16)
we end up with {\em Strassen's algorithm}. The total cost now is $\mathcal{O}(n^{\log_2 7})\sim \mathcal{O}(n^{2.8074})$.


Coppersmith and Winograd improved Strassen's algorithm further, currently sitting at around ${\mathcal O}(n^{2.37})$ operations at the cost of an increased constant.

An important open question is whether we can we construct a fast, practical algorithm\footnote{Although fast matrix multiplication algorithms such as Strassen’s method and the Coppersmith-Winograd algorithm improve the asymptotic complexity of matrix multiplication, they are rarely used in practical numerical linear algebra. Their large constant factors, increased numerical instability, higher memory requirements,  poor cache behavior and high communication overhead in distributed systems typically outweigh the theoretical savings. Consequently, most practical libraries rely on highly optimized implementations of the classical algorithm.} that can multiply two arbitrary $n\times n$ matrices
in ${\mathcal O}(n^{2+\epsilon})$ (or, say, ${\mathcal O}(n^{2}\log n)$) operations.

Yet, in certain cases one may be content with computing $AB$ approximately instead of exactly. By sacrificing precision, this should allow for a reduced complexity algorithm. Randomized algorithms are a natural candidate to realize this goal.

Let $A \in \mathbb{R}^{m \times n}$ and $B \in \mathbb{R}^{n \times p}$. An intuitive approach, popularized by Drineas, Kannan, and Mahoney~\cite{drineas2006fast}, treats the product $AB$ as a sum of $n$ rank-one outer products
$$AB = \sum_{i=1}^n A^{(i)} B_{(i)}.$$
Instead of summing all $n$ terms, we pick $c$ indices $\{i_1, i_2, \dots, i_c\}$ from $\{1, \dots, n\}$ with replacement, based on a probability distribution $\{p_i\}$. We then form the approximation
$$\widetilde{C} = \sum_{t=1}^c \frac{1}{c p_{i_t}} A^{(i_t)} B_{(i_t)}$$
This is equivalent to $AB \approx (AS)(S^T B)$, where $S$ is a sparse sampling and scaling matrix.

The quality of the approximation depends heavily on the sampling probabilities $\{p_i\}$.
If we use  probabilities that are proportional to the Euclidean norms of the columns and rows, 
$$p_i = \frac{\|A^{(i)}\| \|B_{(i)}\|}{\sum_{j=1}^n \|A^{(j)}\| \|B_{(j)}\|},$$
the expected error is bounded by
$$E[\|AB - \widetilde{C}\|_F] \leq \frac{1}{\sqrt{c}} \|A\|_F \|B\|_F.$$
While $\|A\|_F \|B\|_F$ can be large, this bound may be sufficient for applications where the matrices have a  are already noisy or when the dimension $n$ is in the order of millions and $c$ is in the order of 10000, say. However, the cost for computing all the probabilities $p_i$ must be factored into the total cost of matrix-matrix multiplication.


Random projections (like the subsampled Randomized Hadamard Transform or Gaussian sketches) lead to bounds of the form
$$\|AB - (AS)(S^T B)\| \leq \|AB - (AB)_k\| + \epsilon \|A\|_F \|B\|_F,$$
where $(AB)_k$ is the best rank-$k$ approximation of the product, see e.g.~\cite{HMT11,cohen2015optimal}. This indicates that the randomized product effectively captures the most significant singular values of the true product.

\smallskip
In conclusion, randomized sketching can approximate large matrix products efficiently when the matrices share a very large dimension. However, the benefits are limited when the output size is large, when strong spectral guarantees are required, or when the matrices have large stable rank. Consequently, sketching is most effective when the matrix product serves as an intermediate step in a larger algorithm, such as regression, low-rank approximation, or graph computations.





\section*{Exercises}
\addcontentsline{toc}{section}{Exercises}

\begin{myexercise}[\level\level\sep Optimality of the Johnson-Lindenstrauss Lemma]
    \label{prob:optimality_jl}
    Recall that the Johnson-Lindenstrauss lemma states that for any set of $n$ points in $\R^p$ there is a linear map $f\colon \R^p \to \R^d$, which acts almost as an $\eps$-isometry, for some $d \asymp \log n$. We will show that this dependency is indeed optimal.
    
    Fix some $d \in \N$ and $\eps > 0$. A subset $S \subseteq \R^d$ is said to be $\eps$-\textit{separated} if there are no distinct $s, t \in S$ such that $\norm{s-t}_2\leq\eps$. A \textit{packing number} of a subset $T \subseteq \R^d$ is defined as the size of its largest subset that is $\eps$-separated, i.e.
    \begin{equation*}
        \DD(T, \eps) = \max_{\substack{S \subseteq T \\ S \text{ is } \eps-\text{separated}}} \abs{S}.
    \end{equation*}
    \begin{enumerate}[(a)]
        \item Suppose $S$ is an $\eps$-separated subset of $B^d(1)$, the unit ball in $\R^d$. Using a volume argument, prove that
        \begin{equation*}
            \DD\brap{B^d(1), \eps} \leq \brap{\frac{2+\eps}{\eps}}^d.
        \end{equation*}

        \item Consider the set $X = \brac{x_1,\ldots,x_n}$ of $n \geq d$ points in $\R^n$, given by
        \begin{equation*}
            x_i = \begin{cases}
                e_i &\text{if }i < n,\\
                0 &\text{otherwise};
            \end{cases}
        \end{equation*}
        where $e_i$ is $i$-th coordinate vector in $\R^n$. Suppose that $f\colon\R^n\mapsto\R^d$ is an $\eps$-isometry for $X$. Show that 
        \begin{equation*}
            f(X) = \brac{f(x_1), \ldots, f(x_n)}
        \end{equation*}
        is an $(1-\eps)$ separated subset of some translated copy of $B_d(1+\eps)$.

        \item Conclude that
        \begin{equation*}
            n \leq \brap{\frac{3+\eps}{1-\eps}}^d.
        \end{equation*}
    \end{enumerate}
\end{myexercise}

\begin{myexercise}[\level\level\sep Johnson-Lindenstrauss Lemma: Alternative Version]
    \label{prob:jl_lemma}
    The goal of this exercise is to prove the random projection lemma and then to use this result to prove another version of the Johnson-Lindenstrauss lemma.
    \begin{enumerate}[(a)]
        \item Let $P \in \R^{n\times m}$ be the coordinate projection, which maps a vector in $\R^n$ onto its first $m$ coordinates in $\R^m$. Let $z\in S^{n-1}$ be a random vector sampled uniformly on the sphere $S^{n-1}$. Show that 
        \begin{equation*}
            \E\norm{Pz}_2^2 = \frac{m}{n} \norm{z}_2^2.
        \end{equation*}
        
        \item Prove the following statement using the result in Problem~\ref{prob:lipschitz_concentration}: There exists an absolute constant $c > 0$, such that for any $\eps > 0$, the following inequality holds with probability at least $1-2\exp(-c\eps^2 m)$
        \begin{equation}\label{eq:random_proj_lm}
            (1-\eps)\sqrt{\frac{m}{n}}\norm{z}_2\le \norm{Pz}_2\le (1+\eps)\sqrt{\frac{m}{n}}\norm{z}_2.
        \end{equation}
        
        \item Note that the result in (b) is stated for a random vector, while in dimension reduction we wish to find a randomized projection such that it preserves geometry for a fixed set of points. However, it can be shown that the same result \eqref{eq:random_proj_lm} holds when $z\in\R^n$ is a fixed vector, and $P$ is a orthogonal projection onto an $m$-dimensional subspace chosen uniformly at random from all $m$-dimensional subspaces in $\R^n$. In fact, these two models are equivalent.
        
        You can use the mentioned fact \textbf{without proof}. Using (b), show the following result: Let $\mathcal{X}$ be a set of $n$ points in $\R^n$ and let $\eps > 0$. Suppose that 
        \begin{equation*}
            m \geq \frac{C}{\eps^2} \log n.
        \end{equation*}
        
        Consider a random subspace $E$ of dimension $m$ chosen uniformly from all $m$-dimensional subspaces in $\R^n$, and let $P$ be an orthogonal projection on this set. Then with probability at least $1 - 2 \exp(-c \eps^2 m)$, the scaled projection $Q = \sqrt{\tfrac{n}{m}} P$ is an $\eps$-approximate isometry for $\mathcal{X}$, i.e., for all $x, y \in \mathcal{X}$,
        \begin{equation*}
            (1 - \eps) \| x - y \| \le \| Q x - Q y \| \le (1 + \eps) \| x - y\|.
        \end{equation*}
        Here $C,c > 0$ are universal constants.
    \end{enumerate}
\end{myexercise}

\begin{myexercise}[\level\level \sep Random projections flatten vectors]\label{prob:projections_flatten_vectors}
     We saw that random projections are good at preserving euclidean distances. In this exercise we'll study their interaction with the $\ell_\infty$-norm. A vector is considered "flat" if its $\ell_\infty$-norm is small compared to its euclidean norm. We will show that random projections flatten vectors. 
    \begin{enumerate}[(a)]
        \item We are given a deterministic matrix $A \in \R^{m \times n}$, a fixed vector $v \in \R^n$ with $\norm{v} =1$ and a random diagonal matrix $D \in \R^{n \times n}$ that has independent Rademacher entries on its diagonal (so $\mathbb{P}(D_{i,i} = 1) = \mathbb{P}(D_{i,i} = -1) = \frac12$). Prove
        \begin{equation*}
            \mathbb{P}(\norm{ADv}_\infty \geq \norm{A}_\infty t) \leq 2m e^{-t^2/2},
        \end{equation*}
        where $\norm{A}_\infty = \max_{i,j} |A_{i,j}|$.
        \item Let $P \in \R^{m \times n}$ be a uniform random projection matrix. Prove that for some universal constant $c>0$
        $$\mathbb{P} \left( \norm{P}_\infty \geq  \frac{t}{\sqrt{n}} \right) \leq 2mn e^{-ct^2}.$$
        You can use without proof that the distribution of any row of $P$ is a uniform random vector on the euclidean sphere $S^{n-1}$.
        \item Conclude that for $\mathcal{X} \subseteq S^{n-1}$ being a set of unit vectors of cardinality $|\mathcal{X}| = k$, we have that with probability at least $1- \frac Cn$ for all $x \in \mathcal{X}$
        $$ \sqrt{\frac nm }\norm{Px }_\infty \leq C'\sqrt{\frac{\log(nmk)}{m}},$$
        where $C,C'>0$ are sufficiently large universal constants.
    \end{enumerate}
\end{myexercise}

\begin{myexercise}
Implement the Randomized SVD algorithm to approximate the rank-$k$ SVD of a $1000 \times 1000$ Hilbert matrix (a matrix known for rapid singular value decay).\\
(a) Compute the true SVD and the Randomized SVD with k=10 and p=5. \\
(b) Plot the singular values obtained from both methods on the same graph. \\
(c) Report the error.
\end{myexercise}


\begin{myexercise}
Let $A \in \mathbb{R}^{m \times n}$ be a matrix of rank $r$, constructed as
$A = U \Sigma  V^T$, where $U \in \mathbb{R}^{m\times r}$ and $V \in \mathbb{R}^{n\times r}$ have orthonormal columns and $\Sigma =\operatorname{diag}(\sigma_1,\dots,\sigma_r)$. Consider the following two singular value decay regimes:
$$\text{(i) Exponential decay:} \qquad
   \sigma_j = e^{-\alpha j}, \quad \alpha > 0.
$$
$$\text{(ii) Polynomial decay:} \qquad
   \sigma_j = j^{-p}, \quad p > 1.$$
For each regime: \\
(a) Generate a test matrix $A$ of size $1000\times 1000$. \\
(b) Compute a rank-$k$ approximation using randomized SVD without power iterations and
compute the relative Frobenius norm error
$\frac{\|A - A_k^{\mathrm{rSVD}}\|_F}{\|A\|_F}$.
\\
(d) Repeat the experiment with power iterations.\\
(e) Compare the results with the optimal truncated SVD.
\end{myexercise}

\begin{myexercise}
Show that if $S$ is a $(1\pm \eps)$ $\ell_2$-subspace embedding for $A$, then it is also a $(1\pm \eps)$ $\ell_2$-subspace embedding for $U$, where $U$ is an orthonormal basis for the column space of $A$.    
\end{myexercise}

\begin{myexercise}\label{ex:countsketch}
Prove Lemma~\ref{le:countsketchinner}.    
\end{myexercise}

\begin{myexercise}(Comparing sketching matrices in least squares)
Let $A \in \mathbb{R}^{n \times d}$ with $n \gg d$ and let $b \in \mathbb{R}^n$. Consider the overdetermined least-squares problem
$$
x^\star = \arg\min_{x \in \mathbb{R}^d} \|Ax - b\|.
$$

Compare the performance of different sketching matrices $S \in \mathbb{R}^{m \times n}$ for approximating $x^\star$ via the sketch-and-solve approach.
\begin{enumerate}
\item Generate a matrix $A \in \mathbb{R}^{2000 \times 50}$ with i.i.d.\ standard Gaussian entries. Generate a ground-truth vector $x_0 \in \mathbb{R}^{50}$ with i.i.d.\ standard Gaussian entries. Set $b = Ax_0 + \eta$, where $\eta \sim \mathcal{N}(0, 0.01 I)$.
Compute the exact least-squares solution
$x^\star = \arg\min_x \|Ax - b\|$
using a standard solver (QR or SVD).
Record $\|Ax^\star - b\|$ and the runtime.

\item 
For each sketch size $m \in \{100, 200, 400, 800\}$, repeat the following experiment 10 times for each sketching matrix:
\begin{enumerate}
\item Gaussian sketch: $S_{ij} \sim \mathcal{N}(0, 1/m)$.
\item SRHT sketch:
$S = \sqrt{\frac{n}{m}} R H D$, where $D$ is a random sign matrix, $H$ is the Hadamard transform, and $R$ samples $m$ rows uniformly.
\item CountSketch: $S$ has one nonzero per column, with random hash function and random signs.
\end{enumerate}

For each sketch $S$, compute the sketched solution
$\tilde{x} = \arg\min_x \|SAx - Sb\|$.

Record (i) the relative solution error: $\|\tilde{x} - x^\star\| / \|x^\star\|$, (ii) the relative residual error: $\|A\tilde{x} - b\| / \|Ax^\star - b\|)$, and (iii) the runtime for forming $SA$ and solving the reduced problem.
\item 
Plot the average relative residual error versus sketch size $m$ for each sketching method. Plot the average runtime versus sketch size $m$.
\end{enumerate}

Comment on which sketch reaches a given accuracy with the smallest $m$, which sketch is fastest overall, and 
how does sparsity of the sketch affect runtime.

\begin{myexercise}
Repeat the previous exercise with a sparse random matrix $A$, where each row has only 10 nonzero entries (the locations and the entries of the 10 non-zero entries in each row are chosen at random).  How does the performance of CountSketch change relative to Gaussian and SRHT?
\end{myexercise}


\end{myexercise}

\begin{myexercise}
Let $L$ be the graph Laplacian introduced in Definition~\ref{def:graphLaplacian}. 
Let $S \in \mathbb{R}^{m \times n}$ be a sketching matrix (Gaussian, SRHT, or CountSketch).
\begin{itemize}
\item[(a)] 
Define the sketched quadratic form
$$
q_S(x) := \|S L^{1/2} \|^2
$$
and suppose $S$ is an $\ell_2$-subspace embedding for the column space of $L^{1/2} U$.
Show that, with high probability,
$$
   (1-\varepsilon) x^T L x \le q_S(x) \le (1+\varepsilon) x^T L x
   \quad \forall x \in U.
$$

\item[(b)] Generate a random geometric graph with $n = 2000$ nodes  as follows: Sample $n$ points $x_1,\dots,x_n$ independently and uniformly at random in the unit square $[0,1]^2$. Fix a radius $r>0$ and create an undirected edge between nodes $i$ and $j$ if $\|x_i - x_j\| \le r$.
Now, form the graph Laplacian $L$ and
compute the first five nontrivial eigenvectors of $L$ using (i) the full Laplacian, (ii) a sketched Laplacian using a Gaussian sketch, (iii) a sketched Laplacian using CountSketch. For each method, compare runtime, approximation error in eigenvalues,
and subspace error (principal angles between eigenspaces). 
\item[(c)] Why does the presence of the nullspace of $L$ require care when applying sketching?
\end{itemize}

    
\end{myexercise}

